\documentclass{beamer}
\usepackage[utf8]{inputenc}
\usepackage{graphicx}

\newtheorem{definicion}{Definición}
\newtheorem{ejemplo}{Ejemplo}

%%%%%%%%%%%%%%%%%%%%%%%%%%%%%%%%%%%%%%%%%%%%%%%%%%%%%%%%%%%%%%%%%%%%%%%%%%%%%%%
\title[Número &\pi&]{Práctica 11}
\author[Elizabeth]{Elizabeth Hernández Martín}
\date[25-04-2014]{25 de abril de 2014}
%%%%%%%%%%%%%%%%%%%%%%%%%%%%%%%%%%%%%%%%%%%%%%%%%%%%%%%%%%%%%%%%%%%%%%%%%%%%%%%

%\usetheme{Madrid}
%\usetheme{Antibes}
%\usetheme{tree}
%\usetheme{classic}

%%%%%%%%%%%%%%%%%%%%%%%%%%%%%%%%%%%%%%%%%%%%%%%%%%%%%%%%%%%%%%%%%%%%%%%%%%%%%%%
\begin{document}
  
%++++++++++++++++++++++++++++++++++++++++++++++++++++++++++++++++++++++++++++++  
\begin{frame}

  \includegraphics[width=0.15\textwidth]{img/ullesc}
  \hspace*{7.0cm}
  \includegraphics[width=0.16\textwidth]{img/fmatesc}
  \titlepage

  \begin{small}
    \begin{center}
     Facultad de Matemáticas \\
     Universidad de La Laguna
    \end{center}
  \end{small}

\end{frame}
%++++++++++++++++++++++++++++++++++++++++++++++++++++++++++++++++++++++++++++++  

%++++++++++++++++++++++++++++++++++++++++++++++++++++++++++++++++++++++++++++++  
\begin{frame}
  \frametitle{Índice}  
  \tableofcontents[pausesections]
\end{frame}
%++++++++++++++++++++++++++++++++++++++++++++++++++++++++++++++++++++++++++++++  


\section{Primera Sección}


%++++++++++++++++++++++++++++++++++++++++++++++++++++++++++++++++++++++++++++++  
\begin{frame}

\frametitle{Primera Sección}

\begin{definicion}

Pi es la relación entre la longitud una una circuferencia y su diámetro, en geometría euclidiana.
Es un número irracional y una de las constantes matemáticasmás importantes. Se emplea frecuentemente 
en matemáticas, física e ingienería. El valor numérico de $\pi$ aprocimado, trucando a sus primeras cifras, 
es el siguiente:

$\pi$ = 3.14159265358979323846...

\end{definicion}

\end{frame}
%++++++++++++++++++++++++++++++++++++++++++++++++++++++++++++++++++++++++++++++  

\section{Segunda Sección}

%++++++++++++++++++++++++++++++++++++++++++++++++++++++++++++++++++++++++++++++  
\begin{frame}

\frametitle{Segunda Sección}

\begin{block}{Historia de $\pi$}
Pi es el número que se obtiene de dividir la longitud de una circunferencia por su diámetro.
No importa el tamaño de la circunferencia. Grande o pequeña, la proporción
entre su longitud y su diámetro es siempre la misma. Aunque es probable que esta propiedad 
fuera conocidad con anterioridad,los grandes estudiosos de este número fueron los
antiguos griegos, como Anaxágoras, Hipócrates de Quios o Antifronte de Atenas. Anteriormente, 
el valor de $\pi$ se determinaba, casi con seguridad, mediante medidas
experimentales. Arquímedes fue el primero que sepamos, que realizó una estimación teórica de su valor. 
Usando polígonos circunscritos e inscritos, determinó que $\pi$
era mayor que 223/71 y menor que 22/7.

\end{block}

\end{frame}
%++++++++++++++++++++++++++++++++++++++++++++++++++++++++++++++++++++++++++++++  
\section{Fórmulas}
\begin{frame}

\frametitle{Fórmulas}

\begin{block}{Cáculo de $\pi$:}
Por aproximación por integración:
$$\int_{0}^{1} \! \frac{4}{1+x^2}\, dx = 4(atan(1) -atan(0) = \pi $$
Esta integral se puede aproximar numéricamente con una fórmula de cuadratura.

Si se utiliza la regla del punto medio se obtien:

$ \pi \approx \frac{1}{n} \sum\limits_{i=1}^{n}f(x_i)\,$,
 con $f(x) = \frac{4}{1+x^2}\,$,
 $x_i = \frac{i - \frac{1}{2}}{n}$,
 para $i = 1, \dots, n$

\end{block}

\end{frame}

%++++++++++++++++++++++++++++++++++++++++++++++++++++++++++++++++++++++++++++++  

\section{Bibliografía}
%++++++++++++++++++++++++++++++++++++++++++++++++++++++++++++++++++++++++++++++  
\begin{frame}
  \frametitle{Bibliografía}

  \begin{thebibliography}{10}

    \beamertemplatebookbibitems
    \bibitem[Tutorial LaTex]{plan}  
    Tutorial LaTex. 

    \beamertemplatebookbibitems
    \bibitem[Tutorial Beamer]{guia}  
    Turorial Beamer. 

    \beamertemplatebookbibitems
    \bibitem[URL:Wikipedia]{latex} 
    Wikipedia. %\verb|http://www.wikipedia.org/wiki/Número_n|

  \end{thebibliography}
\end{frame}

%++++++++++++++++++++++++++++++++++++++++++++++++++++++++++++++++++++++++++++++  
\end{document}
